\maketitle
\head{本版编排与收录说明}
\noindent\textbf{这是一份按内容顺序整理的中文工作译稿，不是全部已上传材料的完结全译本。}第I、II章及第III章的两个已出版分册正文已经连续收入；第0章当前收入§§1--19。第0章§§20--23、第IV章主体及第V卷预备笔记余文尚未连续译入。未译内容没有用原文扫描页、提要或空白标题代替译文。

\section*{正文次序}
正文按原著内容编号组织：导言；第0章预备知识；第I章概形的语言；第II章若干类态射的初等整体研究；第III章凝聚层的上同调研究；后续待译范围说明；独立编排的第V卷预备笔记已译部分；最后为各原分册的书末资料。第III章第一部分的§§1--5与第二部分的§§6--7直接衔接，不再被第一部分的索引隔开。

第0章原来分载于第一卷、第三卷第一部分、第四卷第一部分。本版将已译的§§1--7、§§8--13及§§14--19归并在一起，但保留 $0_{\mathrm I}$、$0_{\mathrm{III}}$、$0_{\mathrm{IV}}$ 等出处记号。它们标示原出版分册，并非重复的三章。这一安排依内容编号，不声称每节的全部前提均已在前文出现；原著跨章引用仍按原号保留。

\section*{实际收录范围}
\begin{longtable}{@{}p{.25\linewidth}p{.67\linewidth}@{}}
\toprule 材料 & 已收入的连续译文及状态\\\midrule\endhead
第一卷（1960） & 导言；第0章§§1--7；第I章§§1--10，止于（10.15.7）；原书书末资料。承接既有V2校订源码。\\
第二卷（1961） & 第II章§§1--8，止于（8.14.14）；原书参考文献、索引、目录、勘误清单一。书末资料沿用既有合订版依据底本恢复的译排来源。\\
第三卷第一部分（1961） & 第0章§§8--13；第III章§§1--5，止于（5.5.2）；原书书末资料。\\
第三卷第二部分（1963） & 第III章§§6--7，止于（7.9.14），连同证明、记号索引、术语索引、原目录及勘误清单二。\\
第四卷第一部分（1964） & 第0章补篇章首说明及§§14--19。本版新增19.7--19.10，止于（19.10.4）。其中（19.7.1.5）正文使用已核定的完备局部形式，原较弱版本完整保存在核查报告，拓扑范围仍待核。\\
第V卷英文预备笔记 & \texttt{EGAV1-2.ps} 的六页可读材料；\texttt{EGAVaiello.pdf} 的卷首说明及“超平面截面与锥投影”§§1--6。余文不计入已译范围。\\
\bottomrule
\end{longtable}

\section*{译文、勘误与核查的关系}
正文采用已经核实的正确陈述、公式和证明；原文错误写法、修改理由及依据记入脚注。作者正式勘误与本次译者补证分开，较长反例与验证收入独立核查报告。尚未确认的疑点不擅自定为错误；原著本来省略的证明、预备笔记中未写完或无法辨认的内容，不凭通常知识补造成“原文”。本版新增注号以H8开头，旧注号保留。

这一原则用于本轮新增19.7--19.10，并承接第七版对19.3--19.6明确列出的回溯修订。\textbf{它不表示此前全部旧译文中的历史核查意见均已逐项回写正文。}其他既有正文保持原稿；其校订状态须看各条脚注与历史报告。

对作者已经撤销的（19.3.12），正文继续仅保留编号与撤销说明；原陈述及问题在脚注，完整历史陈述和原证明在报告附录。此次19.8.7不再引用该撤销条目来扩展提升结论。作者勘误清单三中与新译范围直接相关的三项亦已纳入书末选译。

第V卷上传载体明确属于预备笔记的翻译和编辑，英译编者为Piotr Blass与Joseph Blass。其重组建议、未完成句子和待补证明不按正式刊行的第I--IV章同等状态标示。与英文合并稿重合的PS载体不重复计作新译文。

\section*{第八版新增范围与核查边界}
本版从19.7小节章首开始，连续处理局部同态的存在性与唯一性、Cohen代数与Cohen环、相对形式光滑性、形式非分歧与形式平展代数，止于（19.10.4）。对应第四卷第一部分分册印刷第104--116页（上传PDF第103--115页）。下一接续点为§20“导子与微分”章首说明及（20.1.1）；§20此次没有计入已译正文。

（19.10.3，i）的满同态断言按一般离散环解释不成立。正文改为准确判据：满同态 $A\to B$ 的核 $I$ 满足 $I=I^2$，当且仅当 $B$ 对 $A$ 形式光滑，亦即形式平展；原句、乘积环反例和双向补证均有明确记录。有关局部同态、有限型模的基环以及交叉引用的作者勘误直接落实在正文，并注明原式。

（19.7.1.5）需特别区分极大理想拓扑与指定理想拓扑。正文明确采用足以覆盖本节实际应用、已给出证明的完备局部形式；不据拓扑适用问题宣布原较弱版本存在反例。原较弱陈述及证明全文保存在报告H8-P01，相关19.7.3首句也单独注明待核。故“连续译至§19末尾”指材料覆盖，不表示该待核推广已经得到独立审定。

\section*{版式与使用}
采用 \texttt{\string\documentclass[UTF8,a4paper]\{ctexbook\}}，由文档类决定默认中文字体、字号、版心、标题及双面分页。不另设页边距，不用自定义页眉。原章号、条目号与公式号保留，默认双面排版所需的空白偶数页不删除。前置部分用罗马数字页码，正文以后连续编号；总目录指向本合订本页面，各原著目录页码仍指原版。

PDF、可独立编译的单文件TeX和分文件工程一并提供。实际编译与检查结果见工程包中的记录。编译、条目覆盖、渲染一致性不代替出版级逐句审校或全书数学证明的独立验证；本轮没有重新审定全部旧译文及历史核查意见。
\tableofcontents
